\section{The Carrion-King — Decay, Renewal, and Transformation}

\subsection{Domain and Epithets}

The Carrion-King is the master of endings that become beginnings. He does not destroy, but transforms—turning death into new life, decay into opportunity, and endings into fresh starts. His followers are harvesters of potential, seeing in every fall the seeds of future growth. The King has no will, only a hunger; no malice, only a rhythm. When something ends, he is there to turn it toward the next beginning.

What crumbles feeds what grows. What dies becomes the soil of tomorrow's triumph. The King does not ask for sacrifice—he simply waits for it.

\subsection{Cultural Origins}

The Carrion-King appears wherever things end and begin again:

\begin{itemize}
\item In \textbf{Acasia}, he is the Rot-King, whose touch turns blighted fields to fertile soil. Farmers leave the first sheaf of grain unharvested for him.
\item In \textbf{Vhasia}, he is the Dust-Accountant. Soldiers leave a button from a fallen comrade on a cairn—the King turns that button into a bullet for the next battle.
\item In \textbf{Aeler}, he is the Under-Gardener. Dwarven crypts have bioluminescent fungus fed by the dead. The King tends the crop.
\item In \textbf{Valewood}, he is the Compost-King, the force that turns fallen leaves to loam. The Lethai do not invoke him. They simply let him work.
\end{itemize}

\subsection{Warrior Archetype: The Grave-Warden}

A Runic Warrior of the Carrion-King is not a hunter of the living. They are a \textit{Grave-Warden}—a warrior who harvests endings and turns them into strength. Their Codex is a compost-stained trowel or a belt of bleached bones. Their Thiasos is a fungal familiar that speaks in spores, a crow that follows battlefields, or a sentient blade forged from the iron of a broken plow.

\subsection{Codex Form}

The Grave-Warden's Codex is made from what was once alive:

\begin{itemize}
\item \textbf{Bone-Chime:} A set of three bleached bones strung on a leather cord. It clicks when death is near. Each bone came from a creature the Runekeeper killed.
\item \textbf{Compost-Stained Trowel:} A small iron trowel, crusted with dark earth. It turns soil—and memory.
\item \textbf{Rot-Cord:} A braided cord of fungus-twined leather. It tightens when something is about to die. It loosens when death is harvested.
\end{itemize}

\subsection{Thiasos Form}

\begin{itemize}
\item \textbf{Fungal Familiar:} A pale fungus that grows from a crack in the Runekeeper's armor. It speaks in spores and can detect decay, rot, or impending death within Near range.
\item \textbf{Battle-Crow:} A black crow with one white feather. It follows battles and funerals. It can carry a message to the dead—though the dead may not answer.
\item \textbf{The Bone-Chime (Sentient):} A set of bones that clicks in the voice of the creature they came from. It can sense when something is about to die. It is never wrong.
\end{itemize}

\subsection{Core Mechanic: Harvest Resonance}

When the Grave-Warden is present when something ends (a life, a battle, a structure, a relationship), they may draw a sliver of that ending's energy. Gain +1 die on your next action, or clear 1 Fatigue. You may use this once per scene.

When the Grave-Warden kills a significant enemy (GM discretion), they may mark 1 Obligation to gain a permanent \textit{Harvest Scar}—a small mark on their body that grants +1 die to a specific type of roll related to that enemy's nature (e.g., +1 die to track wolves after killing a wolf-chieftain).

\subsection{Sample Rites}

\subsubsection{Low Rite: Rite of the Rotting Harvest (5 XP)}

\textbf{Tags:} [DECAY], [HARVEST], [BOOST]

\textbf{Threshold:} Organic matter in early stages of decay, or the remains of a recently ended thing (burnt letter, wilted flower, shattered glass).

Accelerate natural decay to weaken or destroy an organic object (ropes, leather, wood), gaining +2 Effect on Break/Sabotage rolls against it. Alternatively, extract value from an ending: from a defeated enemy, gain +1 die on your next action; from a failed plan, reroll one 1 on your next roll; from a broken item, gain 1 Boon.

\textbf{Obligation:} 1

\subsubsection{Low Rite: Rite of the Borrowed Form (7 XP)}

\textbf{Tags:} [KNOWLEDGE], [SPIRIT], [BOOST]

\textbf{Threshold:} A token from a deceased being (hair, nail, written name).

Channel the essence of what was. Gain one skill die reflecting the deceased's expertise for one scene, or ask the GM one question about their knowledge or abilities.

\textbf{Obligation:} 1

\subsubsection{Standard Rite: Rite of the Fertile Death (8 XP)}

\textbf{Tags:} [TRANSFORM], [AREA], [HEAL], [TERRAIN]

\textbf{Threshold:} Ashes, compost, or the remains of anything that once lived.

Transform death into growth: choose one—create beneficial terrain (cover, concealment, or advantageous positioning) OR grant allies +1 die to healing or recovery rolls for the scene.

\textbf{Obligation:} 2

\subsubsection{Standard Rite: Rite of the Great Consumption (13 XP)}

\textbf{Tags:} [DECAY], [AREA], [SUMMON], [TRANSFORM]

\textbf{Threshold:} A significant amount of organic matter (corpse, fallen tree, collapsed building).

Transform a large area through decay and renewal. Choose two of the following: create difficult terrain that favors you and your allies; summon a Cap 3 swarm of scavengers as temporary allies; generate valuable reagents worth 2 XP (usable for crafting or barter).

\textbf{Obligation:} 2

\subsubsection{High Rite: Rite of the Eternal Cycle (14 XP)}

\textbf{Tags:} [TRANSFORM], [FATE], [SACRIFICE]

\textbf{Threshold:} The complete remains of something significant that has ended (a broken oath, a fallen kingdom's charter, a slain legendary beast).

Complete a transformation cycle. Destroy one major asset, enemy, or obstacle, and create something new of equal or greater value. The player and GM collaborate to define the transformation.

\textbf{Obligation:} 3

\subsection{Patron-Specific Talents}

\subsubsection{Harvest Resonance (4 XP)}

\textbf{Requirement:} The Carrion-King

Once per scene after an enemy is defeated or a significant ending occurs, you may draw a sliver of that ending's energy. Gain +1 die on your next action, or clear 1 Fatigue. You cannot use this talent again until the next scene.

\subsubsection{Rot Sense (6 XP)}

\textbf{Requirement:} The Carrion-King

You can sense decay, rot, or corruption within Near range. You know the general direction and approximate strength (minor, moderate, severe). This is a passive ability.

\subsubsection{Essence Feast (8 XP)}

\textbf{Requirement:} The Carrion-King, Witnessed Scar

When you consume the flesh of a fallen supernatural enemy (spirit, demon, fey), you gain a minor property of theirs for one scene. You must describe the consumption—raw, cooked, or ritual.

\subsection{Roleplaying a Grave-Warden}

\begin{itemize}
\item You do not fear endings. You harvest them. Every death is soil. Every failure is fertilizer.
\item You may be called to clear battlefields, compost blighted fields, or ritually end things that should not continue.
\item You may struggle with the weight of endings. Not everything should end. The King does not distinguish.
\item You may be asked to let something die—a friendship, a hope, a part of yourself. The cycle demands it. The King waits.
\end{itemize}

\subsection{Sample NPC: Warden Kael of the Bone-Chime}

\textbf{Archetype:} Grave-Warden

\textbf{Patron:} The Carrion-King

\textbf{Codex:} A bone-chime that clicks when death is near

\textbf{Thiasos:} A fungal familiar that grows from a crack in his armor

\textbf{Personality:} Quiet, patient, relentless. He has harvested a thousand endings. He does not mourn. He does not celebrate. He works.

\textbf{Conflict:} A village is dying—not from plague or war, but from slow despair. The fields are fallow. The wells are dry. The King wants the ending. Kael is not sure he does.

\subsection{GM Guidance: Intrusions for The Carrion-King}

\begin{itemize}
\item A dead creature the Grave-Warden passed earlier now follows them—not as an enemy, but as a witness. It does not attack. It watches.
\item All food in the party's packs spoils instantly, except for a single loaf of black bread. The bread is warm. It was baked yesterday.
\item A villager the Grave-Warden saved bears a small rot-mark on their hand. It spreads one inch each day.
\item The Grave-Warden's bone-chime clicks at an unexpected moment. Something is about to die. The King is hungry.
\end{itemize}

The thread held. The water carried. That is not a Rite. That is grace.
